\section{公共母线功率平衡}
\label{sec:power-balance}

各子模型通过同一公共交流母线耦合。母线方程要求每小时注入功率与吸收功率相等，以保证不同产品路径之间的能量分配具有统一物理边界。设$P_t^{\mathrm{int}}$为内部关键用电需求，$P_t^{\mathrm{int,ENS}}$为其中未满足的部分，则
\begin{align}
P_t^{\mathrm{src,use}}+P_t^{\mathrm{dis}}+P_t^{\mathrm{H2P}}
={}&P_t^{\mathrm{ch}}+P_t^{\mathrm{EL}}
+P_t^{\mathrm{comp,b}}+P_t^{\mathrm{comp,f}}\notag\\
&+P_t^{\mathrm{cable,s}}
+P_t^{\mathrm{marine}}
+\left(P_t^{\mathrm{int}}-P_t^{\mathrm{int,ENS}}\right).
\label{eq:master-bus}
\end{align}
左侧为进入母线的源侧利用功率、电池放电和氢转电功率；右侧依次为电池充电、电解制氢、基础与弹性算力、海缆外送、海洋用能以及实际得到供电的内部负荷。

电池效率已进入式~\eqref{eq:bess-state}，海缆损耗进入式~\eqref{eq:cable-receive}，PUE进入式~\eqref{eq:dc-it-energy}，因此式~\eqref{eq:master-bus}均采用公共母线计量点处的功率，不再重复扣除转换损耗。源侧同时满足
\begin{equation}
P_t^{\mathrm{src,av}}
=P_t^{\mathrm{src,use}}+P_t^{\mathrm{curt}}.
\label{eq:bus-source-split}
\end{equation}

求解后按式~\eqref{eq:master-bus}重新计算母线输入与输出差，定义残差$r_t^{\mathrm{bus}}$并要求
\begin{equation}
\max_{t\in\mathcal T}|r_t^{\mathrm{bus}}|
\le\varepsilon_{\mathrm{bus}}.
\label{eq:bus-residual}
\end{equation}
$\varepsilon_{\mathrm{bus}}$为数值容差。电池能量递推、氢库存和各输出端口也采用同一时间轴进行守恒复核。
